\section{Bounded rank: the Hamilton--Moitra estimate}\label{sec:bounded}

This section reproduces the radial-isotropic and ordered-coordinate argument of
Hamilton and Moitra~\cite{HM}, with a direct compactness proof of the existence
of the radial scaling (as in~\cite{Orig}). It is used for $d$ below a large
absolute constant, and it also shows that ENP frames exist.

\begin{lemma}\label{lem:HM}
Let $1\le d\le n$, let $X\in\R^{n\times d}$ have $\norm{x_i}^2=a$ for all $i$,
and let $\eta=\opn{X^TX-I}$. There is an ENP frame $W$ with
$\fro{X-W}^2\le\eta d(d-1)$.
\end{lemma}

\begin{proof}
If $d=1$ then $X^TX=na=1$ and $W=X$ works. Let $d\ge2$.

\emph{Step 1: full spark.} Suppose first that every $d\times d$ minor
$\mu_S=\det(X_S)^2$ ($S\subset[n]$, $|S|=d$) is nonzero. Consider
$F(s)=\frac12\log\det(X^T\Diag(e^{2s})X)-a\sum_is_i$. By Cauchy--Binet,
\[
 F(s)=\tfrac12\log\sum_{|S|=d}\mu_S\,e^{T(s)_S},\qquad
 T(s)_S=2\Bigl(\sum_{i\in S}s_i-a\sum_{i}s_i\Bigr),
\]
where $T:\R^n\to\R^{\binom nd}$ is linear and $\sum_ST(s)_S=0$ because each $i$
lies in a fraction $d/n=a$ of the $d$-sets. The function
$\mathcal E(y)=\sum_S\mu_Se^{y_S}$ restricted to the subspace $\im T$ has the
nonempty sublevel set $\{\mathcal E\le\mathcal E(0)\}$, on which
$y_S\le\log(\mathcal E(0)/\mu_S)$ for every $S$ and hence, by $\sum_Sy_S=0$,
also $y_S$ is bounded below. This set is compact, so $\mathcal E|_{\im T}$
attains its minimum at some $T(s^*)$, and $s^*$ minimises $F$.

Let $D=\Diag(e^{s^*})$, $M=(X^TD^2X)^{-1/2}$ and $W=DXM$. The computation in
the proof of \cref{lem:potential} (which does not use $X^TX=I$) gives
$\partial_iF=\norm{w_i}^2-a$, so $W$ is ENP, and since
$w_i=e^{s^*_i}Mx_i$ has norm $\sqrt a$,
\[
 w_i=\sqrt a\,\frac{Mx_i}{\norm{Mx_i}}.
\]

\emph{Step 2: ordered coordinates.} Replacing $X,W$ by $XO,WO$ for a suitable
$O\in O(d)$ changes neither the hypotheses nor the distance, and makes
$M=\Diag(\lambda_1,\dots,\lambda_d)$ with $0<\lambda_1\le\dots\le\lambda_d$.
Then $w_{ij}=\lambda_jx_{ij}/\sqrt{c_i}$ with
$c_i=a^{-1}\sum_\ell\lambda_\ell^2x_{i\ell}^2$, so $w_{ij}$ has the sign of
$x_{ij}$ and $w_{ij}^2-x_{ij}^2=x_{ij}^2(\lambda_j^2/c_i-1)$ is nonpositive and
then nonnegative as $j$ increases. As $\sum_j(x_{ij}^2-w_{ij}^2)=a-a=0$, the
prefix sums $\Delta_i(k)=\sum_{j\le k}(x_{ij}^2-w_{ij}^2)$ are nonnegative,
vanish at $k=d$, and
\[
 \sum_j|x_{ij}^2-w_{ij}^2|=2\max_k\Delta_i(k)\le2\sum_{k=1}^{d-1}\Delta_i(k).
\]
Since $W^TW=I$ and $(X^TX)_{jj}\le1+\eta$,
$\sum_i\Delta_i(k)=\sum_{j\le k}((X^TX)_{jj}-1)\le k\eta$. Equal signs give
$(x_{ij}-w_{ij})^2\le|x_{ij}^2-w_{ij}^2|$, and therefore
\[
 \fro{X-W}^2\le2\eta\sum_{k=1}^{d-1}k=\eta d(d-1).
\]

\emph{Step 3: general $X$.} Let $G$ be a Vandermonde matrix with rows
$(1,\theta_i,\dots,\theta_i^{d-1})$ for distinct $\theta_i$; all its maximal
minors are nonzero. Each maximal minor of $X+hG$ is a polynomial in $h$ with
nonzero leading coefficient, so for all but finitely many $h$ the matrix $X+hG$
is full spark; choose such $h_k\to0$ (so small that no row of $X+h_kG$
vanishes) and normalise each row of $X+h_kG$ to length $\sqrt a$ (a positive diagonal left factor, which preserves full spark).
The resulting $X_k\to X$ have ENP frames $W_k$ with
$\fro{X_k-W_k}^2\le\eta_kd(d-1)$, $\eta_k\to\eta$. The set of ENP frames is
compact; a limit point $W$ of $(W_k)$ is ENP and satisfies the bound.
\end{proof}

\begin{corollary}\label{cor:exist}
For all $1\le d\le n$ an ENP frame exists. Consequently every equal-norm
$X\in\R^{n\times d}$ is within $\fro{X-W}^2\le2\fro X^2+2\fro W^2=4d$ of an ENP
frame $W$.
\end{corollary}

\begin{proof}
Apply \cref{lem:HM} to $x_i=\sqrt a\,e_1$ for all $i$.
\end{proof}
